\paragraph{Phase Kickback with Controlled Hadamard}\label{ex:phase-kickback-with-controlled-hadamard}

In the previous examples, we studied phase kickback using:
\begin{itemize}
    \item The controlled-$X$ (CNOT) gate;
    \item The controlled phase gate $CP(\phi)$.
\end{itemize}
Both examples were relatively straightforward because their eigenstates were easy to identify. In this final example, we use a \definition{Controlled Hadamard Gate}:
\begin{equation}
    CH = \begin{pmatrix}
        I & 0 \\
        0 & H
    \end{pmatrix} =
    \begin{pmatrix}
        1 & 0 & 0 & 0 \\
        0 & 1 & 0 & 0 \\
        0 & 0 & \frac{1}{\sqrt{2}} & \frac{1}{\sqrt{2}} \\
        0 & 0 & \frac{1}{\sqrt{2}} & -\frac{1}{\sqrt{2}}
    \end{pmatrix}
\end{equation}
Recall that the Hadamard gate is:
\begin{equation*}
    H = \frac{1}{\sqrt{2}}\begin{pmatrix}
        1 & 1 \\
        1 & -1
    \end{pmatrix}
\end{equation*}
Unlike the phase gate, its eigenstates are \textbf{not} simply the computational basis states. Instead, the eigenvalues of the Hadamard gate are $\pm 1$, and the corresponding eigenstates are:
\begin{equation*}
    \ket{v_+} = \frac{1}{\sqrt{4-2\sqrt{2}}} \left(\ket{0} + \left(\sqrt{2} - 1\right)\ket{1}\right)
\end{equation*}
\begin{equation*}
    \ket{v_-} = \frac{1}{\sqrt{4+2\sqrt{2}}} \left(\ket{0} - \left(\sqrt{2} + 1\right)\ket{1}\right)
\end{equation*}
To observe phase kickback, we need a \textbf{non-trivial phase}. With positive eigenvalues, the phase is trivial (0), so we will use the negative eigenvalue $-1$ and its corresponding eigenstate $\ket{v_-}$.

\highspace
\begin{examplebox}[: Phase Kickback with Controlled Hadamard]
    Let's prepare the state:
    \begin{itemize}
        \item The \textbf{control qubit} is in the state $\ket{+}$:
        \begin{equation*}
            \ket{+} = \frac{1}{\sqrt{2}}(\ket{0} + \ket{1})
        \end{equation*}
        \item The \textbf{target qubit} is in the eigenstate $\ket{v_-}$ of the Hadamard gate:
        \begin{equation*}
            \ket{v_-} = \frac{1}{\sqrt{4+2\sqrt{2}}} \left(\ket{0} - \left(\sqrt{2} + 1\right)\ket{1}\right)
        \end{equation*}
    \end{itemize}
    Applying the controlled Hadamard gate $CH$ to this state will yield:
    \begin{equation*}
        CH \left(\ket{+} \otimes \ket{v_-}\right) = \frac{1}{\sqrt{2}} \left(\ket{0} \otimes \ket{v_-} + \ket{1} \otimes H\ket{v_-}\right)
    \end{equation*}
    Since $\ket{v_-}$ is an eigenstate of $H$ with eigenvalue $-1$, we have:
    \begin{equation*}
        H\ket{v_-} = -\ket{v_-}
    \end{equation*}
    Therefore, the state after applying $CH$ becomes:
    \begin{equation*}
        \begin{aligned}
            CH \ket{+} \otimes \ket{v_-} &= \frac{1}{\sqrt{2}} \left(\ket{0} \otimes \ket{v_-} + \ket{1} \otimes H \ket{v_-}\right) \\
            &= \frac{1}{\sqrt{2}} \left(\ket{0} \otimes \ket{v_-} + \ket{1} \otimes (-\ket{v_-})\right) \\
            &= \frac{1}{\sqrt{2}} \left(\ket{0} \otimes \ket{v_-} - \ket{1} \otimes \ket{v_-}\right) \\
            &= \frac{1}{\sqrt{2}} \left(\ket{0} - \ket{1}\right) \otimes \ket{v_-} \\
            &= \ket{-} \otimes \ket{v_-}
        \end{aligned}
    \end{equation*}
    That minus sign then appeared as a \textbf{relative phase between the two branches of the control qubit}, which is the essence of phase kickback. The control qubit has been kicked back with a phase of $\pi$ (or $-1$) due to the eigenvalue of the target qubit's state under the Hadamard operation.
\end{examplebox}

\highspace
In general, this example shows that \textbf{phase kickback is not tied to any particular quantum gate}. It is a universal consequence of controlled unitary operators. Whenever the target is prepared in an eigenstate of the controlled unitary, the corresponding eigenvalue is transferred to the control qubit as a relative phase. This is the principle that underlies many quantum algorithms, including the \textbf{Quantum Phase Estimation} algorithm.